\section{Reviewed-argument delivery across a process boundary}
\label{sec:a17-correction-consumption}

This section reports a separately frozen check of the trusted-adapter premise used in
the native-default and native-to-store probes. Those probes assume that a complete
reviewed argument object, supplied per pending tool call, is the object the native
runtime executes. The check asks whether three pinned SDKs deliver each reviewed
object to its intended pending call, and leave the other pending calls unchanged,
when the run is paused for review and then continued, including after a clean process
exit and a new interpreter resumes the persisted state.

The question is deliberately narrow. The application merges the reviewed patch into a
complete five-field argument object and supplies that object for every pending call,
including calls the reviewer did not change. The study therefore measures
transmission fidelity of the reviewed object, not partial-patch inference or
field-level preservation by the SDK. The deterministic model emits the declared calls
and then stops; no language model is called, and no provider request is made. A local
tool body writes its actual runtime call identifier and actual arguments to SQLite,
and a separate SDK-free scorer reads those rows in read-only mode.

The matrix crosses three frameworks, two lifecycles, seven call scenarios, six review
edits and two payload families, giving 504 constructed cells. Lifecycles are an
in-process continuation and an exit after a native pause whose persisted state is
resumed by a newly launched interpreter. Scenarios are one call; two calls with
distinct original arguments; and two calls with equal original arguments, each edited
in the first, second or both positions. Edits are a text replacement, integer zero,
boolean false, empty string, a three-field change and a same-value confirmation.
Payload families are ASCII and text with Unicode, quotes and newlines. The protocol,
case matrix, adapters, scorer and SDK sources were hashed before collection; a
separately recorded twelve-cell development preflight is excluded from these
denominators, and 23 runtime and oracle controls were exercised before the freeze.

\begin{table}[htbp]
\centering\small
\caption{Reviewed-argument delivery: cells whose persisted effect multiset equals the
reviewed expectations exactly. Each lifecycle exercises one native continuation route
per framework; the cells are constructed and are not independent samples.}
\label{tab:a17-correction-consumption}
\begin{tabular}{llrrr}
\toprule
Framework & Lifecycle & Cells & Exact match & Mismatch\\
\midrule
Agno 3.1.10 & in-process continuation & 84 & 24 & 60\\
Agno 3.1.10 & exit and fresh-process resume & 84 & 84 & 0\\
LangChain 1.4.7 / LangGraph 1.2.16 & in-process continuation & 84 & 84 & 0\\
LangChain 1.4.7 / LangGraph 1.2.16 & exit and fresh-process resume & 84 & 84 & 0\\
Pydantic AI 2.60.0 & in-process continuation & 84 & 84 & 0\\
Pydantic AI 2.60.0 & exit and fresh-process resume & 84 & 84 & 0\\
\midrule
Total & & 504 & 444 & 60\\
\bottomrule
\end{tabular}
\end{table}

All 60 mismatching cells belong to one path, the Agno in-process continuation, and
they share one mechanism. In every one of them the first pending call executed with
the complete reviewed arguments supplied for the second pending call, while the second
call executed its own; the fields that differ follow the edit, so the pattern is not
restricted to calls with equal original arguments. Of the 72 multi-call cells on that
path, 60 mismatch; the 24 matching cells are the twelve single-call cells and twelve
cells whose two reviewed payloads coincide, where sharing cannot change the delivered
values. The public issue \#8546 already describes shared human-input state across
pending calls in this framework, so the matrix reproduces a known mechanism and
characterizes its current-version scope. It is not reported as a discovered defect,
not confirmed by the maintainers here, and not counted once per cell.

The in-process arm also differs from the fresh arm in its continuation extension point
and in the object it continues from. A 54-cell ablation in the same frozen family
(Agno only) separates the two. Holding the pause, the reviewed objects and the
recorded arguments fixed, the original-run-object route
(\texttt{continue\_run(run\_response=...)})
mismatched in 12 of 18 cells (all twelve multi-call cells; the six single-call cells matched), while the run-identifier route
(\texttt{continue\_run(run\_id, session\_id, requirements)}) matched in all 18 cells in
the same process and in all 18 cells after a clean exit and a new interpreter. The
mechanism therefore follows the continuation route and the state it continues from, not
the process boundary, and the process contrast alone shows no difference. This is a
two-contrast decomposition rather than a single-factor isolation, because the
entrypoint and object identity change together in the first contrast. The matching
paths are equally bounded: LangChain and Pydantic AI use the same documented
continuation route in both lifecycles, so their arms differ mainly in process and
serialization, and one route per lifecycle is exercised.

The primary observable is the exact multiset of call identifier, tool name and
complete typed payload, compared with the reviewed expectations. Field-level
diagnostics locate the same differences and carry no additional detection power; the
two agreed on every cell. No setup mismatch, unsupported path, collection failure,
worker error, native framework error or blocked network attempt occurred. The
constructed cells are not an ecosystem sample: three frameworks were chosen for prior
support, each lifecycle exercises one route, and no prevalence estimate, confidence
interval or significance test applies. What the check supports is conditional and
descriptive: on the documented routes of two SDKs and on the Agno fresh-process route,
a complete reviewed object reached its intended pending call and the unintended call
was unchanged; on the Agno route that continues from the original in-memory run object
with more than one pending call, the reviewed values of one call were executed by
another, while the run-identifier route was clean in the same process.

\paragraph{Same-identifier substitution.}
A third frozen check (16 cells, Pydantic AI 2.60.0) applies a controlled injection
between the review and the continuation: the state the continuation executes from is
replaced while the native call identifier stays unchanged. When the human decision only
approves the pending call, the continuation
executed the substituted arguments in all four cells that changed the supplied message
history, in both lifecycles; substituting the deferred request object alone changed
nothing (two of two executed the reviewed object), and the two cells that approve with no
substitution also executed the reviewed object. When the decision re-supplies the
reviewed object through the native override argument, the reviewed object was executed
in all eight override cells: six with substitution and two unchanged controls. These three contrasts are conditional on the tested path: matching native call
identifiers did not keep the executed arguments equal to the reviewed ones, while
supplying the complete reviewed object through the native override argument left the
executed arguments unchanged under the same substitution. The framework's
documentation describes client-supplied message history as input outside the boundary
the framework enforces, and does not present native approval as a substitute for
server-side authorization
(\url{https://pydantic.dev/docs/ai/core-concepts/message-history/}, accessed
2026-10-24). The result is therefore an empirical statement about an
application-maintained responsibility boundary, not a framework defect; post-approval
substitution itself is reported in prior work on approval hijacking.

\paragraph{Artifact locations.}
The three supplied archives contain frozen protocols, code, cases, source seals,
raw per-cell databases, scores and reports. Their original freeze SHA-256 values are:
\nolinkurl{Correction-consumption-reproduction.zip},
\nolinkurl{740ccc4fee766849dfaa9fec529c79ba36372cd521138f57672aa586ec89cf8b};
\nolinkurl{Agno-entrypoint-reproduction.zip},
\nolinkurl{f17ade246c0e92371874b9fc6373461f2f42ca1ff85ec9e0d2c8d95dd8076166}; and
\nolinkurl{Reviewed-instance-substitution-reproduction.zip},
\nolinkurl{b5871cf623af9cb81fa5c17ac034c913345c8b5d6ed0836e29ed14da498d7a66}.
In each extracted study directory, \texttt{python verify\_offline.py} checks package,
source and raw-record hashes and recomputes effects from the case definitions using
only the standard library. It needs neither SDK imports nor the recorded installation
paths. Recollection instead requires pinned SDK environments and configured collectors;
their concurrency ladder and no-overwrite/no-replay guards remain unchanged.
\texttt{REPRODUCE-V4.md} documents these separate workflows and the history intervention.
Neither workflow requires model calls.
